\begin{table}[htbp]
\centering
\caption{Standardized Effect Sizes}
\label{tab:sde}
\begin{adjustbox}{max width=\textwidth}
\begin{tabular}{lcccccc}
\hline\hline
Outcome & $\hat{\beta}$ & SE & SD($Y$) & SDE & SE(SDE) & Classification \\
\hline
Education Emp. & $-$0.00102 & 0.00176 & 1.853 & $-$0.0006 & 0.0009 & Null \\
Health Care Emp. & 0.00673 & 0.00115 & 1.313 & 0.0051 & 0.0009 & Small positive \\
Accomm./Food Emp. & 0.00015 & 0.00074 & 1.281 & 0.0001 & 0.0006 & Null \\
Total Private Emp. & 0.00327 & 0.00081 & 1.272 & 0.0026 & 0.0006 & Null \\
Education Hires & $-$0.00353 & 0.00251 & 2.159 & $-$0.0016 & 0.0012 & Null \\
Education Earnings & $-$0.00154 & 0.00107 & 1.853 & $-$0.0008 & 0.0006 & Null \\
\hline\hline
\end{tabular}
\end{adjustbox}
\begin{tablenotes}[flushleft]\small
\item \textit{Notes:} \textbf{Country:} United States. \textbf{Research question:} Does federal regulatory tightening that caused mass for-profit college closures (2013--2018) disrupt county-level labor markets in the education sector and adjacent industries? \textbf{Policy mechanism:} Federal enforcement of the Gainful Employment rule, Cohort Default Rate sanctions, and fraud investigations triggered chain-wide closures of for-profit college systems (Corinthian Colleges, ITT Tech, Education Corporation of America), eliminating over 1,200 campuses as employers across 383 counties. \textbf{Outcome definition:} Log employment (beginning-of-quarter count from QWI), log annual hires, and log average monthly earnings by NAICS sector. \textbf{Treatment:} Continuous---number of for-profit institution closures per county interacted with post-treatment indicator. \textbf{Data:} IPEDS institutional directory (1997--2024) for closure identification; Census QWI (2008--2022) for county $\times$ quarter $\times$ NAICS sector employment outcomes. \textbf{Method:} Two-way fixed effects (county + year) with closure count $\times$ post; Callaway-Sant'Anna (2021) for heterogeneity-robust ATT; chain IV using ITT/Corinthian/ECA exposure. Standard errors clustered at county level. \textbf{Sample:} 383 treated counties (experienced $\geq$1 for-profit closure) and 327 control counties (had for-profit colleges but zero closures), 2008--2022. SDE $= \hat{\beta} / \text{SD}(Y)$ where SD($Y$) is the pre-treatment standard deviation. Classification refers to magnitude, not statistical significance: Large ($|$SDE$| > 0.15$), Moderate ($0.05$--$0.15$), Small ($0.005$--$0.05$), Null ($< 0.005$).
\end{tablenotes}
\end{table}
